\section{Actual content and the logarithmic cancellation bill in Lean}
\label{acr-section}

The ordinary content and residual theorems have finite arithmetic cores
that can be checked directly on integer coordinates and actual prime
factorizations. Their formalization retains the boundary between the
adjugate argument and the oriented factorization argument.

For an integer pair $z=(x,y)$ put
\[
 C(z)=\gcd(|x|,|y|),\qquad
 z_{\rm red}=(x/C(z),y/C(z)),\qquad N(z)=x^2+xy+y^2.
\]
In the Lean definition, $C(z)$ is \texttt{Int.gcd} and the divisions
are integer division. All reconstruction identities below also hold at
the zero pair, under the library's total convention for division by zero.
Primitivity and positivity assertions explicitly exclude the zero pair.

\begin{proposition}[The actual coordinate core]\label{acr-content-core}
For every integer pair $z$,
\[
 C(z)z_{\rm red}=z,\qquad C(z)^2N(z_{\rm red})=N(z).
\]
If $z\ne0$, then $C(z)>0$, $N(z)>0$ and $C(z_{\rm red})=1$.
With the multiplication of $\mathbb Z[\zeta]$, for arbitrary $\tau,W$,
\[
 C(\tau W)^2N((\tau W)_{\rm red})=N(\tau)N(W).
\]
If $C(W)=1$, then $C(\tau W)$, cast to an integer, divides $N(\tau)$.
If also $\tau\ne0$, then $C(\tau W)>0$.
\end{proposition}
\begin{proof}
The actual gcd divides both coordinates, so multiplying their integer
quotients by the gcd reconstructs them, including at zero. Substitute
these two identities into the quadratic norm to obtain its scaling law.
Dividing a nonzero pair by its gcd makes it primitive. Positivity of
the norm follows from $4N(x,y)=(2x+y)^2+3y^2$. Its multiplicativity
gives the product norm identity. Finally the adjugate identities in
Lemma~\ref{pc-content}, combined with the integer gcd's own Bezout
coefficients for the primitive pair $W$, prove the stated divisibility.
The product is nonzero by positivity and norm multiplicativity.
\end{proof}

The module \texttt{ActualProjectiveContent} checks these statements
in nine declarations, including the actual norm-seven example
\[
 (-3+2\zeta)(62-149\zeta)=112+273\zeta=7(16+39\zeta).
\]
It checks both gcds and all displayed norms in this finite instance.
The infinite family of Section~\ref{pc-section} is an ordinary proof;
it is not inferred from this closed instance.

\begin{lemma}[A local reflected depth]\label{acr-local-depth}
Let $n,a,f$ be positive integers and $b,r$ nonnegative integers.
If $|a-nf|=b+nr$, then $n\le a+b$.
\end{lemma}
\begin{proof}
If $a\ge n$, use $b\ge0$. Otherwise $nf\ge n>a$, so the equality
becomes $a+b=n(f-r)>0$. A positive integer multiple of the positive
integer $n$ is at least $n$.
\end{proof}

For a natural number $C$, define the actual prime product
\[
 D(C)=\prod_{p\in\operatorname{primeFactors}(C)}p.
\]
It is positive, with empty product one. For positive $C$ it is the
ordinary radical; the formal total convention also assigns one at zero.

\begin{proposition}[An actual product bill and its strict log threshold]
\label{acr-actual-bill}
Let $n,C,V,V',R,R'$ be positive integers. Suppose that at every prime
$p\mid C$ one has $v_p(V)>0$, $v_p(R)>0$, and
\[
 |v_p(V)-n v_p(R)|=v_p(V')+n v_p(R').
\]
Then $D(C)^n\mid VV'$. If also every prime divisor of $C$ is at least
seven and
\begin{equation}\label{acr-strict-log}
 \log V+\log V'<n\log7,
\end{equation}
then $C=1$.
\end{proposition}
\begin{proof}
Lemma~\ref{acr-local-depth} gives
$n\le v_p(V)+v_p(V')=v_p(VV')$. Thus $p^n\mid VV'$ at every
prime in the actual support of $C$. Distinct prime powers are coprime;
their product is $D(C)^n$, proving the bill.

Strict monotonicity of the logarithm on positive reals and the finite
product and power laws give the exact equivalence
\[
 \log V+\log V'<n\log7\quad\Longleftrightarrow\quad VV'<7^n.
\]
If $C>1$, choose an actual prime divisor $p$ of $C$. Its $n$th power
divides $VV'$, so $7^n\le p^n\le VV'$, a contradiction.
\end{proof}

The local depth premise is supplied ordinarily by the oriented
calculation of Lemma~\ref{il-content}. It has not been silently
discharged from the integer gcd theorem in Lean. The bill is for the
product of distinct prime divisors, not for $C$ itself. The strict
inequality cannot be weakened to a non-strict one: the norm-seven
family has $C=7$, $V=7$, $V'=7^{n-1}$ and equality.

\paragraph{Connecting the global radical to a positive integer.}
With the actual prime-logarithm definitions of Section~\ref{apl-section},
the finite product formula proves, for every natural $C$,
\[
 R_0(C)=\log D(C).
\]
Consequently the identity and height transfer of Section~\ref{ags-section}
become
\[
 \mathcal S_0(N)=\log N-3\log D(N),\qquad
 t-(\delta+b+l)/3\le\log D(N),
\]
with exactly the same explicit height, signed-tail and small-mass
hypotheses. This verifies the passage from finite prime sums to the
usual integer radical without adding any analytic premise.

\paragraph{Exact formal scope.}
The modules \texttt{ResidualReflection\allowbreak Arithmetic} and
\texttt{ActualResidual\allowbreak LogThreshold} contain seven and six proved
declarations respectively. They use integer absolute values,
\texttt{Nat.factorization}, actual finite prime products and
\texttt{Real.log}. Their source-level audit checks the explicit
local reflection premise, positivity, the prime lower bound and
strictness of the threshold. The total natural-number statements
occasionally allow zero where the mathematical application uses
positive inputs; this does not omit any required positive-domain premise.

Together with the nine content declarations and the nine declarations
in \texttt{ActualGlobal\allowbreak SignedBridge}, the joint fresh run checks
31 new declarations and 53 unchanged local dependency declarations.
It uses Lean 4.32.0, warnings as errors, and the pinned Mathlib revision
\texttt{81a5d257c8e410db227a6665ed08f64fea08e997}.
The compiled Mathlib cache is reused. All 84 declarations have explicit
axiom queries; the union is only \texttt{propext},
\texttt{Classical.choice} and \texttt{Quot.sound}.
The ordinary proofs precede these implementations.

The Eisenstein UFD, general integral inverse, rational-point existence,
analytic mean estimates and arbitrary-root signed bounds are not
formalized by these modules. Independent review and exact finite
replay complement the scoped compilation; none constitutes a proof
or disproof of $abc$.
